\documentclass{article}
\usepackage[T1]{fontenc}
\usepackage{lmodern}
\usepackage{graphicx}
\usepackage{amsmath,amssymb}
\usepackage[a4paper,margin=2cm]{geometry}
\usepackage[absolute,overlay]{textpos}
\setlength{\TPHorizModule}{1mm}
\setlength{\TPVertModule}{1mm}
\usepackage[indonesian]{babel}
\usepackage{hyperref}
\setlength{\emergencystretch}{2em}
% unit: Halaman 8

\begin{document}

% === Halaman 8 (python/native) ===

• Solusi dari superincreasing knapsack (yaitu b1, b2, …, bn) mudah dicari 

sebagai berikut (berarti sama dengan mendekripsikan cipherteks 

menjadi plainteks semula): 

1.

Jumlahkan semua bobot di dalam barisan.

2.

Bandingkan bobot total dengan bobot terbesar di dalam barisan. Jika 

bobot terbesar lebih kecil atau sama dengan bobot total, maka ia 

dimasukkan ke dalam knapsack, jika tidak, maka ia tidak dimasukkan.

3.

Kurangi bobot total dengan bobot yang telah dimasukkan, kemudian 

bandingkan bobot total sekarang dengan bobot terbesar selanjutnya. 

Demikian seterusnya sampai seluruh bobot di dalam barisan selesai 

dibandingkan.

4.

Jika bobot total menjadi nol, maka terdapat solusi persoalan 

superincreasing knapsack , tetapi jika tidak nol, maka tidak ada solusinya. 

8

\end{document}
